% !TeX program = xelatex
\documentclass[12pt,a4paper,oneside]{report}
% !TeX root = ../main.tex
\XeTeXgenerateactualtext=1
\usepackage[margin=23mm,top=24mm,bottom=24mm,headheight=27pt]{geometry}
\usepackage{amsmath,amssymb,mathtools,bm}
\usepackage{booktabs,array,longtable,tabularx}
\usepackage[table]{xcolor}
\usepackage{tikz}
\usetikzlibrary{arrows.meta,calc,patterns}
\usepackage[most]{tcolorbox}
\usepackage{titlesec,needspace}
\usepackage{hyperref}
\usepackage{bookmark}
\usepackage{xepersian}

% !TeX root = ../main.tex
\settextfont[Scale=1.06,BoldFont=Amiri-Bold.ttf]{Amiri-Regular.ttf}
\setlatintextfont[BoldFont=texgyretermes-bold.otf,ItalicFont=texgyretermes-italic.otf,BoldItalicFont=texgyretermes-bolditalic.otf]{texgyretermes-regular.otf}

% !TeX root = ../main.tex
\definecolor{navy}{HTML}{163747}
\definecolor{teal}{HTML}{127C80}
\definecolor{gold}{HTML}{B18542}
\definecolor{paper}{HTML}{F3F7F7}
\definecolor{muted}{HTML}{5E7178}
\definecolor{alert}{HTML}{A45538}

% مشخصات جزوه؛ عنوان و نام‌ها فقط در این فایل تغییر می‌کنند.
\newcommand{\BookTitle}{داده‌کاوی}
\newcommand{\BookSubtitle}{از دادهٔ خام تا الگو و مدل قابل اعتماد}
\newcommand{\CourseInstructor}{دکتر شیرین شعاعی}
\newcommand{\BookCompilers}{احسان شهبازی}

\hypersetup{unicode=true,colorlinks=true,linkcolor=navy,urlcolor=teal,
 pdftitle={داده‌کاوی ـ کتاب آموزشی و تعاملی},pdfauthor={احسان شهبازی},pdfsubject={پیش‌پردازش، طبقه‌بندی، رگرسیون، خوشه‌بندی و تحلیل وابستگی}}

% !TeX root = ../main.tex
\setlength{\parindent}{0pt}
\setlength{\parskip}{5pt plus 1pt}
\linespread{1.12}
\raggedbottom
\setcounter{tocdepth}{1}
\setcounter{secnumdepth}{2}
\allowdisplaybreaks[2]
\emergencystretch=2em
% Page styles use explicit boxes to keep bidi direction groups balanced.
\makeatletter
\def\ps@lecture{%
  \def\@oddhead{\vbox{\hbox to\textwidth{\small
    \rl{\BookTitle}\hfil\rl{\leftmark}}%
    \kern4pt\hrule height .35pt}}%
  \let\@evenhead\@oddhead
  \def\@oddfoot{\hfil\thepage\hfil}%
  \let\@evenfoot\@oddfoot}
\makeatother
\pagestyle{lecture}
\renewcommand{\chaptermark}[1]{\markboth{#1}{}}
\titleformat{\chapter}[display]{\color{navy}\bfseries\Huge}{\large فصل \thechapter}{8pt}{}[\vspace{4pt}\color{gold}\titlerule]
\titlespacing*{\chapter}{0pt}{0pt}{18pt}
\titleformat{\section}[block]{\Large\bfseries\color{navy}}{\thesection}{.6em}{}
\titleformat{\subsection}[block]{\large\bfseries\color{teal}}{\thesubsection}{.6em}{}

% !TeX root = ../main.tex
\tcbset{enhanced,breakable,boxrule=.45pt,arc=1.5mm,
 left=3mm,right=3mm,top=2mm,bottom=2mm,before skip=9pt,after skip=9pt,
 fonttitle=\bfseries,coltitle=white,colback=paper,colframe=teal}
\newtcolorbox{definition}[1]{breakable=false,title={تعریف | #1},colframe=navy,colback=navy!3!white}
\newtcolorbox{result}[1]{breakable=false,title={نتیجهٔ کلیدی | #1},colframe=teal,colback=teal!4!white}
\newtcolorbox{important}[1]{breakable=false,title={دقت کنید | #1},colframe=gold,colback=gold!6!white}
\newcounter{ex}[chapter]
\newenvironment{example}[1]{\refstepcounter{ex}\begin{tcolorbox}[breakable=false,title={مثال \thechapter.\arabic{ex} | #1},colframe=muted,colback=white]}{\end{tcolorbox}}

% Structured exercises: the web builder turns answers into closed disclosures.
\newenvironment{selftests}{\begin{enumerate}\setlength{\itemsep}{9pt}}{\end{enumerate}}
\newenvironment{selftest}[1]{\item\textbf{#1}\par}{\par}
\newcommand{\answer}[1]{\par{\color{teal}\bfseries پاسخ و توضیح: }#1\par}

% !TeX root = ../main.tex
\newcommand{\edit}[1]{\par\begingroup\small\color{alert}\textbf{تکمیل یا اصلاح علمی.} #1\par\endgroup}
\newcommand{\source}[1]{\par{\small\color{muted} منبع مطالعه: #1}\par}
\newcommand{\E}{\mathbb E}
\newcommand{\Var}{\operatorname{Var}}
\newcommand{\Cov}{\operatorname{Cov}}
\newcommand{\MSE}{\operatorname{MSE}}
\newcommand{\Bias}{\operatorname{Bias}}
\newcommand{\SE}{\operatorname{SE}}
\newcommand{\iid}{\overset{\mathrm{iid}}{\sim}}
\newcommand{\N}{\mathcal N}
\newcommand{\R}{\mathbb R}
\newcommand{\Prob}{\mathbb P}
\newcommand{\sol}{\par\textbf{حل.} }
\newcommand{\figcap}[1]{\par\smallskip{\small\color{muted}#1}\par}

\newcommand{\labref}[1]{\par\smallskip{\small\color{teal}\textbf{آزمایشگاه پیشنهادی.} #1}\par}
\newcommand{\checkanswer}[1]{\par\textbf{پاسخ کوتاه.} #1\par}

\providecommand{\tightlist}{\setlength{\itemsep}{0pt}\setlength{\parskip}{0pt}}
\providecommand{\passthrough}[1]{#1}



\hypersetup{pdftitle={داده‌کاوی؛ از مسئله تا کشف الگو},pdfauthor={دکتر شیرین شعاعی؛ احسان شهبازی}}
\begin{document}
\setcounter{chapter}{0}
\chapter{داده‌کاوی؛ از مسئله تا کشف الگو}
\begin{center}{\small\color{muted} استاد درس: دکتر شیرین شعاعی\\[2mm] گردآورنده جزوه: احسان شهبازی}\end{center}

\source{اسلاید مقدمه، صفحات PDF ۱۹ تا ۵۸؛ کتاب، فصل ۱، صفحات چاپی ۲۴ تا ۳۳ و فصل ۶، بخش ۶٫۳.}

\section{داده، الگو و تصمیم چه تفاوتی دارند؟}
ثبت شدن هزاران خرید، به‌تنهایی توضیح نمی‌دهد که چرا مشتریان خرید را ترک می‌کنند. یک جدول بزرگ فقط مجموعه‌ای از مشاهده‌هاست. وقتی رابطه‌ای تکرارشونده میان مشاهده‌ها پیدا می‌کنیم، به یک الگو می‌رسیم؛ وقتی آن الگو را در زمینهٔ مسئله تفسیر می‌کنیم، می‌تواند مبنای تصمیم شود. داده‌کاوی \lr{Data Mining} فرایند کشف الگوهای معتبر، تازه، بالقوه مفید و قابل فهم در داده است. هر یک از این چهار صفت مهم است: رابطه‌ای که تنها حاصل تصادف باشد معتبر نیست؛ رابطهٔ بدیهی ممکن است اطلاعات تازه ندهد؛ رابطهٔ دقیق ولی بی‌ربط به تصمیم، فایدهٔ عملی کمی دارد؛ و خروجی نامفهوم را به‌سختی می‌توان ارزیابی یا به‌کار گرفت.

\begin{definition}{بازیابی با یادگیری یکسان نیست}
پرسش «مشتریان با خرید بیش از یک میلیون را پیدا کن» یک شرط از پیش معلوم را روی داده اجرا می‌کند. پرسش «با رفتار ماه گذشته، چه کسی احتمالاً ماه آینده خرید می‌کند؟» نیازمند یادگیری رابطه‌ای از داده و تعمیم آن به موارد ندیده است. اجرای سریع یک پرس‌وجوی پایگاه داده، به‌خودی خود داده‌کاوی نیست.
\end{definition}

کتاب داده‌کاوی را بخشی از فرایند گسترده‌تر کشف دانش از پایگاه داده، یا \lr{KDD}، معرفی می‌کند. دادهٔ خام ممکن است ناقص، سوگیرانه یا نامتناسب با الگوریتم باشد. حتی پس از ساخت مدل، باید اعتبار، قابلیت تفسیر و اثر استفاده از آن بررسی شود. بنابراین کیفیت نتیجه حاصل کل زنجیره است، نه صرفاً انتخاب یک الگوریتم مشهور.

\begin{example}{سه پرسش دربارهٔ یک فروشگاه}
«چه کالاهایی دیروز فروخته شدند؟» توصیف مستقیم داده است. «چه کالاهایی معمولاً با هم خریده می‌شوند؟» تحلیل وابستگی است. «درآمد هفتهٔ آینده چقدر است؟» رگرسیون یا پیش‌بینی سری زمانی است، بسته به نحوهٔ صورت‌بندی و وابستگی زمانی. یک جدول واحد می‌تواند ورودی چند مسئله باشد؛ نوع مسئله را پرسش و خروجی مورد نیاز تعیین می‌کند.
\end{example}

\section{چرخهٔ کشف دانش و نقش هر مرحله}
چرخه با فهم مسئله آغاز می‌شود: تصمیم‌گیرنده کیست، خروجی چه زمانی لازم است، خطا چه هزینه‌ای دارد و هنگام پیش‌بینی چه اطلاعاتی واقعاً در دسترس خواهد بود؟ سپس داده انتخاب و بررسی اولیه می‌شود. تعداد سطرها، نوع ویژگی‌ها، مقدارهای گمشده، توزیع هدف و منشأ داده را بررسی می‌کنیم. این مرحله را کاوش یا پروفایل‌سازی داده می‌نامند.

در پیش‌پردازش، خطاها و تکرارها را بررسی، نوع داده را درست و نمایش مناسب را تولید می‌کنیم. در مدل‌سازی، پارامترهای یک روش را با دادهٔ آموزش برآورد می‌کنیم. در ارزیابی، می‌پرسیم آیا مدل روی دادهٔ ندیده و با معیار متناسب با هدف کار می‌کند؟ در تفسیر، بررسی می‌کنیم چه چیزی آموخته شده و چه فرضی لازم است. در استقرار، مدل در یک فرایند واقعی استفاده و عملکرد آن در زمان پایش می‌شود. چارچوب \lr{CRISP-DM} نیز بر فهم کسب‌وکار، فهم داده، آماده‌سازی، مدل‌سازی، ارزیابی و استقرار تأکید دارد.

\begin{result}{بازگشت به عقب بخشی از کار است}
چرخهٔ داده‌کاوی یک خط مستقیم نیست. اگر مدل روی یک زیرگروه ضعیف باشد، ممکن است به جمع‌آوری دادهٔ بیشتر، اصلاح تعریف هدف یا نمایش بهتر ویژگی‌ها نیاز داشته باشیم. افزایش پیچیدگی مدل همیشه نخستین پاسخ مناسب نیست.
\end{result}

فرض کنید هدف، تشخیص ایمیل ناخواسته است. باید پیش از دیدن مدل روشن کنیم که ایمیل متعلق به زمان ارسال است یا پاسخ کاربر پس از دریافت هم در ویژگی‌ها آمده است. متن ایمیل را به ویژگی عددی تبدیل می‌کنیم، مدل را روی ایمیل‌های گذشته آموزش می‌دهیم و آن را روی ایمیل‌های مستقل می‌سنجیم. اگر واژگان فرستندگان تغییر کنند، عملکرد آینده می‌تواند افت کند؛ پایش و بازآموزی ادامهٔ چرخه‌اند.

\section{چهار وظیفهٔ اصلی و معنای برچسب}
در یادگیری نظارت‌شده \lr{Supervised Learning}، برای مثال‌های آموزشی، پاسخ مطلوب داریم. در طبقه‌بندی \lr{Classification}، پاسخ از مجموعه‌ای از طبقه‌هاست؛ مثلاً «هرزنامه/عادی». در رگرسیون \lr{Regression}، پاسخ عددی است؛ مثلاً قیمت یک خانه. هدف مشترک این است که رابطهٔ ورودی و پاسخ را برای نمونهٔ جدید بیاموزیم، نه اینکه پاسخ‌های آموزش را حفظ کنیم.

در خوشه‌بندی \lr{Clustering}، برچسب گروه‌ها از قبل به الگوریتم داده نشده است. الگوریتم بر اساس یک تعریف شباهت، ساختاری در داده پیدا می‌کند. گروه‌های حاصل لزوماً با طبقه‌بندی انسانی یکی نیستند. در تحلیل وابستگی \lr{Association Analysis}، به دنبال اقلام یا رخدادهایی هستیم که با هم مشاهده می‌شوند؛ معمولاً خروجی به شکل مجموعهٔ پرتکرار یا قاعده است.

\begin{center}
\begin{tabularx}{\dimexpr\textwidth-5pt\relax}{p{24mm}XX}
\toprule
وظیفه & خروجی & نمونهٔ پرسش\\
\midrule
طبقه‌بندی & برچسب گسسته & آیا این پیام ناخواسته است؟\\
رگرسیون & مقدار عددی & هزینهٔ سفر چقدر خواهد شد؟\\
خوشه‌بندی & گروه‌بندی بدون برچسب هدف & مشتریان چه تیپ‌های رفتاری دارند؟\\
وابستگی & مجموعه یا قاعدهٔ هم‌وقوعی & کدام کالاها با هم خریداری می‌شوند؟\\
\bottomrule
\end{tabularx}
\end{center}

توصیفی بودن با بی‌نیازی از ارزیابی برابر نیست. خوشه‌هایی که با تغییر کوچک مقیاس از بین می‌روند یا قواعدی که فقط در چند خرید ظاهر شده‌اند، ممکن است قابل اتکا نباشند. همچنین «پیش‌بینی» به معنای کشف رابطهٔ علّی نیست: رابطه‌ای که برای پیش‌بینی مفید است الزاماً نشان نمی‌دهد با مداخله روی یک متغیر، دیگری تغییر خواهد کرد.

\section{نمایش داده، ویژگی و هدف}
معمولاً هر سطر یک شیء یا مشاهده و هر ستون یک ویژگی \lr{Feature} است. برای دادهٔ جدولی می‌نویسیم
\[X=(x_{ij})\in\mathbb R^{n\times d},\qquad \mathbf y=(y_1,\ldots,y_n).\]
این نمایش عددی به آن معنا نیست که دادهٔ اولیه حتماً عددی بوده است؛ واژه‌ها، رنگ‌ها و تاریخ‌ها باید با رعایت معنایشان به نمایش قابل استفاده تبدیل شوند. «شناسهٔ دانشجو» ممکن است عدد باشد ولی فاصلهٔ دو شناسه معنای آموزشی ندارد. در مقابل، فاصلهٔ دو اندازه‌گیری قد معمولاً معنا دارد.

کتاب میان دادهٔ رکوردی، تراکنشی، سند، گراف و دادهٔ مرتب مانند سری زمانی تفاوت می‌گذارد. برای تراکنش، هر سطر یک مجموعه از اقلام است. در سند، هر مؤلفه می‌تواند وزن یک واژه باشد. در گراف، رابطهٔ میان اشیا نیز بخشی از داده است. اگر داده زمانی باشد، جابه‌جا کردن تصادفی سطرها می‌تواند مسئلهٔ واقعی را عوض کند؛ هنگام طراحی آزمون باید ترتیب و گروه‌بندی رعایت شود.

پیش از ساخت مدل، واحد مشاهده را دقیق تعریف کنید. اگر هر بیمار چند ردیف داشته باشد، تقسیم ردیف‌ها میان آموزش و آزمون ممکن است بیمار یکسان را در هر دو قرار دهد. در این حالت باید تقسیم بر اساس بیمار انجام شود. این نکته از همان ابتدا بر معنی «نمونهٔ جدید» اثر می‌گذارد.

\section{نخستین الگوریتم: نزدیک‌ترین همسایه‌ها}
ایدهٔ \lr{k-Nearest Neighbors} یا \lr{k-NN} این است که نمونهٔ جدید احتمالاً به نمونه‌های مشابه خود شبیه است. برای پیش‌بینی، فاصلهٔ آن تا داده‌های آموزش را حساب می‌کنیم، $k$ فاصلهٔ کوچک‌تر را برمی‌گزینیم و در طبقه‌بندی رأی اکثریت می‌گیریم:
\[\widehat y(\mathbf x)=\underset{c}{\operatorname{argmax}}\sum_{i\in N_k(\mathbf x)}\mathbf1\{y_i=c\}.\]
برای دادهٔ عددی، یک انتخاب معمول فاصلهٔ اقلیدسی است:
\[d(\mathbf x,\mathbf z)=\sqrt{\sum_{j=1}^d(x_j-z_j)^2}.\]
این فاصله یک انتخاب مدل‌سازی است و باید با معنای ویژگی‌ها سازگار باشد. در رأی وزن‌دار، همسایهٔ نزدیک‌تر سهم بیشتری دارد؛ اگر وزن معکوس فاصله انتخاب شود باید تکلیف فاصلهٔ صفر را جداگانه روشن کرد.

\begin{example}{رأی‌گیری و نقش مقیاس}
پنج ایستگاه نزدیک، به‌ترتیب برچسب‌های «آفتابی، ابری، ابری، آفتابی، ابری» دارند. برای $k=5$ رأی اکثریت «ابری» است. اگر فقط نزدیک‌ترین ایستگاه را بگیریم، پاسخ «آفتابی» می‌شود. تغییر $k$ می‌تواند پاسخ را تغییر دهد؛ انتخاب آن باید با اعتبارسنجی انجام شود. اگر فاصله را از سن و درآمد بسازیم، اختلاف چند هزار واحد درآمد می‌تواند اختلاف چند سال سن را تقریباً حذف کند. بنابراین همسایه‌ها تابع نمایش و مقیاس داده‌اند.
\end{example}

$k$ کوچک، مرز تصمیم جزئی و حساس به نویز می‌سازد؛ $k$ بسیار بزرگ می‌تواند تفاوت محلی طبقه‌ها را محو کند. عدد فرد برای مسئلهٔ دوطبقه‌ای احتمال تساوی رأی را کم می‌کند، ولی جایگزین قاعدهٔ روشن برای شکستن تساوی نیست. آموزش این روش عمدتاً ذخیرهٔ داده است، اما پیش‌بینی سادهٔ یک مشاهده حدود $O(nd)$ هزینه دارد. ساختمان‌های جست‌وجو در ابعاد مناسب می‌توانند کمک کنند، ولی در ابعاد زیاد مفهوم نزدیکی دشوارتر می‌شود.

\labref{در آزمایشگاه نزدیک‌ترین همسایه‌ها، $k$، محل نقطهٔ پرسش و نویز داده را تغییر دهید. پیش از دیدن پاسخ، برچسب را حدس بزنید و بعد بررسی کنید کدام همسایه‌ها رأی داده‌اند.}

\section{چالش‌های داده‌کاوی و پرسش‌های پایانی}
بزرگی داده فقط تعداد سطرها نیست. تعداد زیاد ویژگی‌ها، سرعت ورود، نوع‌های متفاوت و توزیع داده میان چند منبع نیز مهم‌اند. زیاد شدن ویژگی‌ها فضای جست‌وجو را بزرگ و نقاط را در فضا پراکنده می‌کند. دادهٔ فراوان ولی نامرتبط یا سوگیرانه لزوماً مدل خوبی نمی‌سازد. داده‌کاوی از آمار برای سنجش عدم قطعیت، از یادگیری ماشین برای تعمیم، از بهینه‌سازی برای برازش و از پایگاه داده برای ذخیره و پردازش کارآمد استفاده می‌کند.

در گزارش هر پروژه، این زنجیره را قابل پیگیری بنویسید: پرسش، واحد مشاهده، دادهٔ قابل دسترس هنگام پیش‌بینی، نمایش ویژگی، تقسیم داده، مدل پایه، معیار ارزیابی و نتیجهٔ آزمون مستقل. اگر نتوانید معیار موفقیت را پیش از دیدن نتیجه بیان کنید، مقایسهٔ روش‌ها مبهم می‌ماند.
\section{از فاصلهٔ عددی تا پیش‌بینی؛ یک اجرای دستی}
فرض کنید نقطهٔ پرسش $q=(2,1)$ و چهار مشاهدهٔ آموزش به‌شکل زیر باشند:
\begin{center}
\begin{tabular}{ccc}
\toprule مشاهده & مختصات و برچسب & فاصله تا پرسش\\\midrule
$A$ & $(1,1),\ y=0$ & $1$\\
$B$ & $(2,3),\ y=1$ & $2$\\
$C$ & $(4,1),\ y=1$ & $2$\\
$D$ & $(0,0),\ y=0$ & $\sqrt5$\\\bottomrule
\end{tabular}
\end{center}
برای نمونه، فاصلهٔ $D$ برابر $\sqrt{(2-0)^2+(1-0)^2}=\sqrt5$ است. پس با یک همسایه، برچسب صفر و با سه همسایه، برچسب یک پیش‌بینی می‌شود. برابری فاصلهٔ $B,C$ در این پرسش رأی سه‌تایی را مبهم نمی‌کند، چون هر دو انتخاب می‌شوند؛ اگر مرز انتخاب از میان هم‌فاصله‌ها بگذرد باید قرارداد ثابتی داشته باشیم. فاصله فقط می‌گوید کدام مشاهده نزدیک‌تر است؛ برچسب با رأی‌گیری به دست می‌آید. این دو مرحله را از هم جدا کنید.

برای استفادهٔ واقعی، ابتدا ویژگی‌ها و تعریف فاصله را تعیین و تبدیل لازم را فقط از آموزش یاد می‌گیریم؛ سپس $k$ را با اعتبارسنجی انتخاب می‌کنیم. نقطهٔ پرسش نباید در محاسبهٔ میانگین و مقیاس آموزش شرکت کند. پس از پیش‌بینی نیز لازم است بدانیم همسایه‌ها از چه گروهی آمده‌اند: نزدیکی عددی، به‌تنهایی دلیل علّی یا تضمین کیفیت نیست.
\labref{در آزمایشگاه، نقطهٔ پرسش را در مرز دو رنگ بگذارید؛ پاسخ یک و چند همسایه را مقایسه کنید و نام همسایه‌های رأی‌دهنده را از جدول بررسی کنید.}

\section{خودآزمایی و تمرین}
پاسخ را ابتدا با بیان خودتان بنویسید؛ سپس توضیح را مقایسه کنید.
\begin{selftests}
\begin{selftest}{نوع مسئله را از روی خروجی تعیین کنید}
گروه‌بندی مشتریان بدون برچسب، پیش‌بینی مبلغ خرید و پیش‌بینی خرید/عدم خرید چه تفاوتی دارند؟
\answer{به‌ترتیب خوشه‌بندی، رگرسیون و طبقه‌بندی‌اند. در دو مورد آخر پاسخ آموزشی داریم، ولی نوع خروجی متفاوت است. یکسان بودن جدول ورودی، وظیفه را یکسان نمی‌کند.}
\end{selftest}
\begin{selftest}{بازیابی یا یادگیری؟}
آیا پیدا کردن دانشجویی با یک شمارهٔ مشخص، داده‌کاوی است؟
\answer{بازیابی مستقیم است؛ شرط را از قبل می‌دانیم و رابطه‌ای برای تعمیم به مشاهدهٔ ندیده نمی‌آموزیم. پیش‌بینی خطر افت تحصیلی از ویژگی‌های قبلی، مسئلهٔ یادگیری متفاوتی است.}
\end{selftest}
\begin{selftest}{رأی همسایه‌ها را حساب کنید}
در مثال این فصل، خروجی برای یک و سه همسایه چیست؟ اگر هر چهار همسایه رأی دهند چه مسئله‌ای پیش می‌آید؟
\answer{برای یک همسایه صفر و برای سه همسایه یک است. با چهار همسایه رأی‌ها برابر می‌شوند؛ باید قاعدهٔ شکستن تساوی را از قبل تعیین کرد. عوض کردن آن پس از دیدن پاسخ درست، ارزیابی را مخدوش می‌کند.}
\end{selftest}
\begin{selftest}{چرا آزمون نهایی محل انتخاب نیست؟}
بیست مقدار $k$ را روی آزمون می‌سنجیم و بهترین را گزارش می‌کنیم. مشکل چیست؟
\answer{نتیجهٔ آزمون وارد انتخاب مدل شده است. بهترین امتیاز از میان تلاش‌های متعدد معمولاً خوش‌بینانه است؛ تنظیم را با اعتبارسنجی انتخاب و فقط نتیجهٔ نهایی را روی آزمون مستقل گزارش می‌کنیم.}
\end{selftest}
\begin{selftest}{نشت پاسخ را تشخیص دهید}
برای پیش‌بینی قیمت خانه، ستون قیمت فروش را هم در ورودی قرار می‌دهیم. دقت عالی چه معنایی دارد؟
\answer{پاسخ از قبل در اختیار مدل است؛ این نتیجه کیفیت پیش‌بینی واقعی را نشان نمی‌دهد. باید فقط اطلاعاتی را وارد کرد که در لحظهٔ پیش‌بینی در دسترس‌اند. همین بررسی دربارهٔ ویژگی‌های پس از رخداد نیز لازم است.}
\end{selftest}
\begin{selftest}{همراهی چه زمانی علت نیست؟}
خریداران یک کالا بیشتر کالای دیگری می‌خرند. آیا تخفیف کالای اول حتماً فروش دومی را افزایش می‌دهد؟
\answer{خیر؛ ویژگی مشتری، فصل یا عامل مشترک می‌تواند هر دو خرید را توضیح دهد. الگوی مشاهده‌ای برای پیش‌بینی مفید است، اما اثر مداخله به طراحی علّی، آزمایش یا فرض‌های اضافی نیاز دارد.}
\end{selftest}
\end{selftests}

% !TeX root = ../main.tex

\end{document}
